\section{Discussion}
\label{sec:discussion}

For running AI weather models on a modest research setup with a single consumer-grade GPU, our results
support running the network in fp16 with the physical state,
normalisation and denormalisation in fp32. The change in skill is
statistically detectable with 52 paired dates but small in practical
terms -- less than an hour of lead time up to day 8 -- whereas the
savings in time and energy are a factor of two to four for four of
the five models, and about 1.4 for SFNO. Two conditions
come with this recommendation. First, the conversion must leave the
state in fp32: the one-line conversion of a whole model does not, and
for some models it fails silently, as for SFNO, whose complex weights
lose their imaginary parts. Second, fp16's numeric range
must be checked for each model, as we did with hooks on every module
(\S\ref{sec:results-headroom}); the margins we found differ by more than
a factor of ten between models, and a model trained or fine-tuned
differently may have larger activations.

Across the five models, which span 62 million to 1.26
billion parameters (Table~\ref{tab:models}), the sensitivity to fp16
follows neither model size nor fp32 skill. The two smallest models,
FastNet and FourCastNet, depart furthest from their fp32 forecasts
(Figure~\ref{fig:fp16-departure}), but Pangu-Weather and SFNO, of
intermediate size, depart least, and Aurora, the largest, lies in
between. SFNO has among the largest z500 errors of the five models yet is the
least affected by fp16, while Pangu-Weather and Aurora, the most skilful,
are affected little and moderately. The one model clearly more affected by fp16, FourCastNet, does not stand
out by size or skill: it has the largest changes in RMSE (up to 1.4\%,
against about 0.5\% at most for the others), the largest departure from
its fp32 forecasts relative to its error at day 1 (8\% in the mean over variables, against
1--5\%) and a tropical z500 penalty of up to 12\%
(\S\ref{sec:results-fp16}, \S\ref{sec:results-where}). A plausible cause is its
architecture, in particular its adaptive Fourier filters on the flat
latitude--longitude grid, but we have not isolated it; SFNO, which also
works in spectral space but on the sphere, is the least sensitive of the
five. With five models, this is an observation rather than a
test: the models also differ in step length (Pangu-Weather rounds 10
times in a 10-day forecast, the others 40 times), in which operations
always stay in fp32 (SFNO's spherical harmonic transforms), in
resolution, and in how fast they amplify small perturbations, all of
which affect how far the fp16 forecasts drift.

The computing power used for machine learning is increasingly
concentrated in industry, and academic groups have become less
represented in compute-intensive research
\citep{ahmed2020dedemocratization, besiroglu2024compute}. Hardware with
more memory has also become harder to afford: retail prices of consumer GPUs
have risen sharply, and most for the models with the most memory
\citep{techtimes2026gpu}, as graphics memory is increasingly taken up by AI
hardware \citep{trendforce2026memory}. This study also shows what direct access to
a modest local GPU offers over rented or remote hardware. Short tests
between experiments, such as those that located a memory leak in SFNO's
software or FastNet's sensitive layer in int8, took minutes and could
follow one another without waiting in a job queue; exclusive use of the GPU made timings reproducible
and allowed the board power to be sampled every 50\,ms, which shared or
virtualised GPUs often do not permit; the input data and the software
environments of five models stayed in place between experiments; and
the GPU cost nothing while it was idle. Together, these considerations
make it worthwhile to make AI weather models, efficient as they already
are, more efficient still: in fp16, every model studied here runs on a
single 16\,GB consumer GPU, Aurora without streaming its weights from CPU
memory, 1.3 to 3.5 times faster than in fp32 and at a fraction of the
energy, with practically no loss of skill.

The results are subject to several limitations. They cover five
models, chosen because they can be run in PyTorch on a 16\,GB GPU; larger
models, and models such as GraphCast that run in other frameworks, may
behave differently. All forecasts start in 2024, a single year, and are
verified against ERA5 rather than observations, for a few variables at
daily lead times up to 10 days. The storms are four cases and the
ensemble one model on 13 dates; they show that fp16 does not visibly
degrade these forecasts, not that it never does. The costs were measured
on one consumer GPU with one version of PyTorch, and both the speed-up
and the size of the rounding errors depend on the hardware and on the
kernels the library selects. Finally, the sensitivity experiments
measure departures from the fp32 forecast over the first day, on 13
dates; they rank the parts of each network but do not show how errors
introduced there affect skill at longer lead times.

The recommendation assumes hardware with fast fp16 arithmetic
(\S\ref{sec:quantisation}). Google's TPUs, on which several AI
weather models are developed, multiply matrices in bf16
\citep{wang2019bfloat16}; there bf16 is the natural 16-bit format, and
our results suggest it costs more skill (\S\ref{sec:results-bf16}). The
sensitivity experiments indicate where to look for a remedy: for
Pangu-Weather, Aurora and FastNet, the first or last components cause
most of the departure (\S\ref{sec:results-sensitivity}), so keeping these
in fp32 should remove much of it; for FourCastNet, whose sensitivity is
spread over the whole network, it may not.

Simple 8-bit quantisation, with one scale per output channel,
costs far more skill than fp16 (\S\ref{sec:results-8bit}). Newer formats
share a scale among small blocks of values \citep{rouhani2023microscaling},
and current GPUs, including the consumer GPU used here, compute in
floating-point formats down to 4 bits \citep{nvidia2025rtx50}. Block
scaling addresses the failure we found in FastNet, where a single input
with large weights set the int8 steps for most rows of a layer
(\S\ref{sec:results-8bit}), because each value then shares its scale only
with a few neighbours. Whether such formats, perhaps with quantisation-aware
fine-tuning, make 8- or 4-bit inference of weather models as accurate as
fp16 is open, as is whether they would be faster for models such as
Aurora, whose memory and time are dominated by intermediate results
rather than weights. The paired
design and module hooks used here would answer the first question
directly.
